\section{Testing and approximation lower bounds}

The lower comparison in \cref{obj:thm:uniform-annotation-frontier} combines rare-label testing with a many-cell construction based on polynomial approximation. The first component isolates the information supplied by complete records about rare-arm outcomes. The second connects a reciprocal approximation gap to priors on complete-record tables, with a common distribution of treatment--covariate marginals. We first fix the testing distances and signed-measure norm used in these arguments.



The probability distance measures separation between observation laws, whereas the signed-measure norm fixes the normalization of an approximation certificate. The next three results supply, respectively, the finite-alphabet Pinsker inequality \citep[Lemma 2.5]{Tsybakov2009Nonparametric}, Markov's polynomial inequality \citep[Theorem 1.1]{Pierzchala2016Polynomial}, and the finite alternation representation used in polynomial moment matching \citep[Remark 4 and Appendix E, identity (34)]{WuYang2016Entropy}. These are standard tools; the new statistical construction begins with their overlap-scaled specialization below.

\begingroup\footnotesize
\begin{lemmav}{lem:pinsker-finite}[Pinsker inequality on finite alphabets]
Let \(\nu,\nu'\) be probability measures on a finite set. Define their total variation distance and natural-logarithm Kullback--Leibler divergence by
\[
\operatorname{TV}(\nu,\nu')
=\frac12\sum_z|\nu(z)-\nu'(z)|,
\qquad
D(\nu\Vert\nu')
=
\begin{cases}
\displaystyle\sum_z \nu(z)\log\frac{\nu(z)}{\nu'(z)},
& \nu'(z)=0\Longrightarrow \nu(z)=0\ \text{for every }z,\\
+\infty,& \text{otherwise}.
\end{cases}
\]
Each summand with \(\nu(z)=0\) is defined to be zero. Then, with the inequality interpreted in the extended nonnegative reals,
\[
\operatorname{TV}(\nu,\nu')
\le \sqrt{\frac{D(\nu\Vert\nu')}{2}}.
\]
\end{lemmav}
\endgroup

\Cref{obj:lem:pinsker-finite} expresses testing separation in terms of the natural-logarithm divergence. For the approximation component, the relevant control instead concerns how rapidly a bounded polynomial can vary over its interval.

\begin{lemmav}{lem:markov-polynomial}[Markov inequality for polynomials]
Let \(L\) be a nonnegative integer, \(g\) a real polynomial of degree at most \(L\), and \(M_b\in\mathbb R\). Suppose that
\[
|g(x)|\le M_b \qquad \text{for every } x\in[-1,1].
\]
Then
\[
|g'(x)|\le L^2 M_b \qquad \text{for every } x\in[-1,1].
\]
\end{lemmav}

The degree-squared derivative bound identifies the interval scale relevant to the reciprocal function below. A complementary duality result represents its approximation error by a signed measure that annihilates all moments through the prescribed degree.

\begingroup\footnotesize
\begin{lemmav}{lem:finite-polynomial-duality}[Finite polynomial approximation duality]
Let \(v_-<v_+\), let \(g:[v_-,v_+]\to\mathbb R\) be continuous, and let \(L\ge0\) be an integer. There exist nodes
\[
v_-\le v_1<\cdots<v_{L+2}\le v_+
\]
and real weights \(w_1,\ldots,w_{L+2}\) such that
\[
\sum_{i=1}^{L+2}|w_i|=1,
\qquad
\sum_{i=1}^{L+2}w_i v_i^h=0
\quad\text{for every integer }0\le h\le L,
\]
and
\[
\sum_{i=1}^{L+2}w_i g(v_i)
=
\inf_{Q:\,\deg Q\le L}
\sup_{v\in[v_-,v_+]}|g(v)-Q(v)|,
\]
where the infimum ranges over real polynomials of degree at most \(L\). Equivalently, the finite signed measure
\[
\sigma=\sum_{i=1}^{L+2}w_i\delta_{v_i}
\]
has support on at most \(L+2\) points and satisfies
\[
\|\sigma\|_{\mathrm{TV}}=1,
\qquad
\int v^h\,d\sigma(v)=0
\quad\text{for every integer }0\le h\le L,
\qquad
\int g\,d\sigma
=
\inf_{Q:\,\deg Q\le L}
\sup_{v\in[v_-,v_+]}|g(v)-Q(v)|.
\]
Here \(\delta_{v_i}\) is the unit point mass at \(v_i\), and \(\|\sigma\|_{\mathrm{TV}}=\sum_{i=1}^{L+2}|w_i|\) is the total mass of the variation measure of \(\sigma\).
\end{lemmav}
\endgroup

The finite support in \cref{obj:lem:finite-polynomial-duality} makes the certificate suitable for a finite latent distribution. The following specialization retains a fixed reciprocal gap while placing its scale at the interval length divided by the squared degree.

\begin{lemmav}{lem:shrinking-cone-dual}[Finite reciprocal-gap certificate]
Let the parameters satisfy:
\begin{itemize}
\item \textbf{(Degree.)} \(L\) is an integer with \(L\ge 2\).
\item \textbf{(Interval.)} \(B>0\).
\item \textbf{(Overlap.)} \(0<\epsilon\le 1/4\).
\end{itemize}
Set
\[
\alpha=\frac{B}{100L^2}.
\]
There exist strictly increasing nodes \(0\le v_0<\cdots<v_{L+1}\le B\) and real weights \(c_0,\ldots,c_{L+1}\) defining the finite signed measure
\[
\sigma=\sum_{i=0}^{L+1}c_i\delta_{v_i},
\]
where \(\delta_{v_i}\) is the unit point mass at \(v_i\), such that
\[
\|\sigma\|_{\mathrm{TV}}
=\sum_{i=0}^{L+1}|c_i|=1,
\qquad
\int v^h\,d\sigma(v)=0
\quad\text{for every integer }0\le h\le L,
\]
and
\[
\int\frac{\alpha}{\alpha+(1-\epsilon)v}\,d\sigma(v)
\ge\frac18.
\]
Here \(\|\sigma\|_{\mathrm{TV}}\) denotes the total mass of the variation measure of \(\sigma\).
\end{lemmav}

The certificate separates two mathematical roles: its vanishing moments support the observation comparison, and its reciprocal integral supplies the target separation. Both properties hold on the same finite support under the degree, interval, and overlap conditions in \cref{obj:lem:shrinking-cone-dual}.

The testing priors use a two-cell construction for the rare-label component and a certificate-based construction for the many-cell component. We first specify the two-cell tables and the variation measure needed to turn signed certificate weights into latent probabilities.

\begingroup\small
\begin{definitionv}{synth_9}[Rare-label probability tables]
For \(d\ge2\), \(0<\epsilon\le1/4\), and \(0<\delta\le1/8\), define \(P_1\) and \(P_0\) on \(\{1,\ldots,d\}\times\{0,1\}^2\) by
\[
\begin{aligned}
P_1(X=j,A=1,Y=1)&=\frac{\epsilon}{2}\left(\frac12+\delta\right),&
P_1(X=j,A=1,Y=0)&=\frac{\epsilon}{2}\left(\frac12-\delta\right),\\
P_0(X=j,A=1,Y=1)&=\frac{\epsilon}{2}\left(\frac12-\delta\right),&
P_0(X=j,A=1,Y=0)&=\frac{\epsilon}{2}\left(\frac12+\delta\right),
\end{aligned}
\qquad j\in\{1,2\}.
\]
For both tables, set
\[
P_1(X=j,A=0,Y=y)=P_0(X=j,A=0,Y=y)=\frac{1-\epsilon}{4},
\qquad j\in\{1,2\},\quad y\in\{0,1\},
\]
and assign zero probability to every record with \(j>2\).
Thus \(P_1\) is the positive-sign law and \(P_0\) the negative-sign law of \(\cref{obj:thm:rare-label-floor}\). In the rare-label branch of \(\cref{obj:def:prior-handle}\), use these tables with
\[
\delta=\frac{1}{16\sqrt{\max\{1,n\epsilon\}}},
\qquad
\Pi_1=\delta_{P_1},\qquad \Pi_0=\delta_{P_0},
\]
where \(\delta_P\) denotes the Dirac probability measure at the table \(P\).
\end{definitionv}
\endgroup

These tables have the same treatment--covariate marginal \leanref{sym:P_{XA}}{\ensuremath{P_{XA}}}, while their treated outcome means carry opposite perturbations. They supply the rare-label comparison in \cref{obj:thm:rare-label-floor}, including the benchmark with the exact marginal table supplied. For the many-cell construction, the variation measure assigns nonnegative mass to the certificate's atoms and leaves their signs available for the outcome assignment.

\begin{definitionv}{synth_6}[Total-variation measure]
For the selected finite signed measure \(\sigma\), let \(|\sigma|\) denote its total-variation measure. For every measurable set \(E\),
\[
|\sigma|(E)=\sup_{\{E_1,\ldots,E_r\}}\sum_{i=1}^{r}|\sigma(E_i)|,
\]
where the supremum ranges over all finite measurable partitions of \(E\). In particular, if \(\sigma=\sum_{i=1}^{r}c_i\delta_{v_i}\) with distinct atoms \(v_i\), then
\[
|\sigma|=\sum_{i=1}^{r}|c_i|\delta_{v_i}.
\]
\end{definitionv}

The construction below chooses its branch using the public information scales. In the affine branch, a common latent distribution determines the arm masses, a core branch carries the signed outcome assignment, and a reservoir supplies the remaining raw mass before normalization.

\begingroup\scriptsize\setlength{\arraycolsep}{2pt}%
\advance\hsize by 8pt\advance\linewidth by 8pt
\begin{algorithmv}{def:prior-handle}[Two testing prior construction]
The testing priors \(\Pi_1,\Pi_0\) are measures on complete-record probability tables, defined for public parameters \(n,m,d\in\mathbb N\) and \(\epsilon\in\mathbb R\) satisfying
\begin{itemize}
\item \textbf{(Complete-record budget.)} \(n\ge1\).
\item \textbf{(Alphabet size.)} \(d\ge2\).
\item \textbf{(Overlap range.)} \(0<\epsilon\le1/4\).
\end{itemize}
Their construction is as follows.
\begin{enumerate}
\item Put
\[
S=n\epsilon,\qquad
\ell=\log(\exp(1)+S),\qquad
N=n+m,\qquad
x=\frac d{N\epsilon\ell}.
\]
If \(S<\exp(4096)\) or \(x^2\le S^{-1}\), set
\[
\delta=\frac1{16\sqrt{\max\{1,S\}}}.
\]
For \(r\in\{0,1\}\), define the two-cell table by
\[
P_r(X=j,A=a,Y=y)=
\begin{cases}
\displaystyle
\frac12\epsilon^a(1-\epsilon)^{1-a}
\left(\frac12+(2r-1)a\delta\right)^y
\left(\frac12-(2r-1)a\delta\right)^{1-y},
&j\in\{1,2\},\\[4pt]
0,&3\le j\le d,
\end{cases}
\qquad a,y\in\{0,1\}.
\]
Output
\[
\Pi_1=\delta_{P_1},\qquad \Pi_0=\delta_{P_0}.
\]
Otherwise continue with the following steps.

\item Set
\[
\begin{aligned}
u&=128n,\qquad w=128N,\qquad M=u+w,\\
L&=\lceil8\ell\rceil,\qquad
B=\frac L{1000M},\qquad
\alpha=\frac B{100L^2},\\
b_0&=\frac\alpha\epsilon,\qquad
K_*=\min\left\{d-1,\left\lfloor\frac1{100b_0}\right\rfloor\right\}.
\end{aligned}
\]

\item Fix a witness supplied by \cref{obj:lem:shrinking-cone-dual}, represented as
\[
\sigma=\sum_{i=1}^{L+2}\omega_i\delta_{v_i},
\]
with defining properties
\[
\begin{gathered}
0\le v_1<\cdots<v_{L+2}\le B,\qquad
\sum_{i=1}^{L+2}|\omega_i|=1,\\
\sum_{i=1}^{L+2}\omega_i v_i^h=0
\quad(h=0,\ldots,L),\\
\sum_{i=1}^{L+2}\omega_i
\frac{\alpha}{\alpha+(1-\epsilon)v_i}\ge\frac18.
\end{gathered}
\]
At each atom of positive variation mass, define the polar sign
\[
h(v)=\frac{\sigma(\{v\})}{|\sigma|(\{v\})}\in\{-1,1\}.
\]

\item Define the raw covariate and arm masses by
\[
w_*(v)=b_0+v,\qquad
S_1(v)=\alpha+(1-\epsilon)v,\qquad
S_0(v)=(1-\epsilon)b_0+\epsilon v,
\qquad 0\le v\le B.
\]
For each rare cell \(i=1,\ldots,K_*\), independently draw a latent coordinate \(V_i\) and a branch flag \(c_i\in\{0,1\}\), with \(c_i=1\) denoting the core branch, according to
\[
\begin{aligned}
\Pr(V_i\in dv,\ c_i=1)
&=\frac{\alpha}{S_1(v)}|\sigma|(dv),\\
\Pr(V_i=0,\ c_i=0)
&=1-\int\frac{\alpha}{S_1(v)}|\sigma|(dv).
\end{aligned}
\]
The filler and core branches remain distinct when the core measure has an atom at zero. For \(r\in\{0,1\}\), set
\[
\mu_1^{(r)}(v,c)=
\begin{cases}
\bigl(1+(2r-1)h(v)\bigr)/2,&c=1,\\
0,&c=0,
\end{cases}
\qquad
\mu_0^{(r)}(v,c)=0.
\]
Both priors use raw arm masses \(S_1(v),S_0(v)\).

\item Let \(V\) have the common latent-coordinate marginal just defined, and put
\[
\bar w=\mathbb E[w_*(V)],\qquad
p_*=1-K_*\bar w,\qquad
Q=p_*+\sum_{i=1}^{K_*}w_*(V_i).
\]
Append a reservoir cell with raw covariate mass \(p_*\), propensity \(1/2\), and both outcome means zero. Set the remaining cells null.

\item For \(r\in\{0,1\}\), define the normalized complete-record table by
\[
P_r(X=j,A=a,Y=y)=
\begin{cases}
\displaystyle
\frac{S_a(V_j)}Q
\left[
y\mu_a^{(r)}(V_j,c_j)
+(1-y)\bigl(1-\mu_a^{(r)}(V_j,c_j)\bigr)
\right],
&1\le j\le K_*,\\[6pt]
\displaystyle\frac{p_*}{2Q}\mathbf1\{y=0\},
&j=K_*+1,\\[6pt]
0,&K_*+1<j\le d,
\end{cases}
\qquad a,y\in\{0,1\}.
\]
Output the pushforward priors
\[
\Pi_1=\operatorname{Law}(P_1),\qquad
\Pi_0=\operatorname{Law}(P_0)
\]
of the common finite product latent distribution under the respective table maps.
\end{enumerate}
\end{algorithmv}
\endgroup

The testing priors \leanref{sym:\Pi_a}{\ensuremath{\Pi_1,\Pi_0}} in \cref{obj:def:prior-handle} are distributions over normalized probability tables. In the affine branch, the reservoir mass \(p_*\) complements the expected total rare-cell mass, while \(Q\) is the realized total raw mass. Normalizing by \(Q\) therefore retains the variation in rare-cell masses across latent draws. The core and filler flags specify distinct outcome assignments even when they share a latent coordinate at zero.

The reciprocal certificate supplies the separation ingredient for the covariate-weighted target in \cref{obj:def:target}. Its moment conditions supply the approximation ingredient for comparing observations under the two hypotheses. The next result assembles these roles under the large-information and affine-regime conditions, including the overlap validity of every normalized table.

\begin{lemmav}{lem:normalized-affine-testing}[Normalized affine testing priors]
There exists a numerical constant \(c>0\) such that the following holds for every \(n,m,d\in\mathbb N\) and \(\epsilon\in\mathbb R\). Put
\[
S=n\epsilon,\qquad N=n+m,\qquad
\ell=\log(\exp(1)+S),\qquad x=\frac{d}{N\epsilon\ell},
\]
where \(S\) is the rare-arm information scale, \(N\) is the total marginal-information sample size, \(\ell\) is the regularized logarithmic scale, and \(x\) is the normalized many-cell difficulty. Suppose that:
\begin{itemize}
\item \textbf{(Public indices.)} \(n\ge1\) and \(d\ge2\).
\item \textbf{(Overlap parameter.)} \(0<\epsilon\le1/4\).
\item \textbf{(Information scale.)} \(S\ge\exp(4096)\).
\item \textbf{(Affine regime.)} \(x^2>S^{-1}\).
\end{itemize}
For every signed certificate \(\sigma\) satisfying the certificate conditions at the affine tuning parameters in \cref{obj:def:prior-handle}, and every latent configuration in that construction, the normalized laws under both hypotheses belong to \(\mathcal M_{d,\epsilon}\) of \cref{obj:def:model}, and their treatment--covariate marginal tables \(P_{XA}\) agree. For each such \(\sigma\), the two product priors obtained by pushing the common product latent distribution through the normalized-law construction also induce exactly the same distribution of \(P_{XA}\).

The selected priors \(\Pi_1\) and \(\Pi_0\) in \cref{obj:def:prior-handle} induce exactly the same distribution of \(P_{XA}\), and
\[
R(n,m,d,\epsilon)\ge c\min\{1,x^2\}.
\]
Here \(R(n,m,d,\epsilon)\) is the minimax squared-error risk in \cref{obj:def:risk}, with the infimum over all clipped Borel randomized rules in the original fixed-size two-channel experiment.
\end{lemmav}

For a common latent configuration, the two hypotheses have identical normalized marginal tables because their arm masses and normalization agree. Pushing the common latent distribution through the two table maps consequently gives the same distribution of \(P_{XA}\) under both priors. Across latent configurations, the marginal table itself can vary. Thus the prior property concerns equality of distributions over marginal tables, together with equality under the specified common-latent coupling.

The risk conclusion in \cref{obj:lem:normalized-affine-testing} concerns the joint observation law of complete and auxiliary records in \cref{obj:def:risk}. Its observation comparison incorporates the complete-record outcomes and their association with the latent marginal table. The common marginal distribution describes the auxiliary channel's unconditional law under the priors; the stated lower bound accounts for both channels together. The resulting many-cell scale depends on total marginal information through \(N\), while the logarithmic degree calibration depends on rare-arm information through \(S\).

Finally, \cref{obj:thm:rare-label-floor} supplies the capped inverse-\(S\) lower bound throughout the public parameter range. It controls the branch \(x^2\le S^{-1}\) and, with the fixed numerical threshold absorbed into the comparison constant, the bounded-information branch \(S<\exp(4096)\). In the remaining branch, \cref{obj:lem:normalized-affine-testing} supplies the capped \(x^2\) lower bound. Together these comparisons give the lower order
\[
\min\{1,S^{-1}+x^2\}
\]
specified by \cref{obj:def:rate-handle}. Combined with the attaining-estimator upper comparison in \cref{obj:thm:uniform-annotation-frontier}, they establish the uniform minimax comparison for the original two-channel experiment.
